\documentclass[british]{article}

\usepackage[margin=1.5in]{geometry}

\title{Relative duality in complex analytic geometry}
\author{J.P. Ramis, G. Ruget, J.L. Verdier}
\date{1971}

% Standard
\usepackage{amssymb,amsmath,amscd}
\usepackage{enumerate}
\usepackage{tikz-cd}
\usepackage{mathtools}
\usepackage{booktabs}


%Bibliography
\usepackage[backend=biber,maxnames=99,maxalphanames=4]{biblatex}
\addbibresource{IM-13-1971-261.bib}


% Colours
\usepackage{xcolor}


% Fonts
\usepackage{mathrsfs}
\usepackage{fouriernc}


% Things for Quarto/pandoc
\usepackage{calc,array}
\usepackage{longtable}
\providecommand{\tightlist}{%
\setlength{\itemsep}{0pt}\setlength{\parskip}{0pt}}
\newcounter{none}
\usepackage{graphicx}
\makeatletter
\newsavebox\pandoc@box
\newcommand*\pandocbounded[1]{% scales image to fit in text height/width
  \sbox\pandoc@box{#1}%
  \Gscale@div\@tempa{\textheight}{\dimexpr\ht\pandoc@box+\dp\pandoc@box\relax}%
  \Gscale@div\@tempb{\linewidth}{\wd\pandoc@box}%
  \ifdim\@tempb\p@<\@tempa\p@\let\@tempa\@tempb\fi% select the smaller of both
  \ifdim\@tempa\p@<\p@\scalebox{\@tempa}{\usebox\pandoc@box}%
  \else\usebox{\pandoc@box}%
  \fi%
}
% Set default figure placement to htbp
\def\fps@figure{htbp}
\makeatother
\usepackage{kvsetkeys}


% Header and footer
\usepackage{fancyhdr}
\usepackage{lastpage}
\usepackage{datetime2}
\pagestyle{fancy}
\fancypagestyle{plain}{}
\fancyhf{}
\renewcommand{\headrulewidth}{0pt}
\cfoot{\small\thepage\ of \pageref*{LastPage}}
\lfoot{\footnotesize Build date: \today}


% Theorem environments
\usepackage{amsthm}
\newenvironment{itenv}[1]
  {\phantomsection\par\smallskip\noindent\textbf{#1.}\itshape}
  {\par\smallskip}
\newenvironment{rmenv}[1]
  {\phantomsection\par\smallskip\noindent\textbf{#1.}\rmfamily}
  {\par\smallskip}


% Shortcuts
\newcommand{\oldpage}[1]{\marginpar{\footnotesize$\Big\vert$ \textit{p.~#1}}}


%% Document %%

% hyperref last, as is "good practice"

\usepackage{hyperref}
\hypersetup{colorlinks,linkcolor={purple!50!black},citecolor={purple!50!black},urlcolor={purple!80!black}}

\begin{document}

\maketitle
\thispagestyle{fancy}

\setcounter{footnote}{0}


%% Content %%

\emph{This document is a translation into English of the following:}

\begin{quote}
Ramis, J.P., Ruget, G., and Verdier, J.L. ``Dualité relative en
géométrie analytique complexe''. \emph{Inventiones Math.} \textbf{13}
(1971), 261--283.
\href{https://eudml.org/doc/142096}{eudml.org/doc/142096}
\end{quote}

\emph{The translator (\href{https://thosgood.net}{Tim Hosgood}) takes
full responsibility for any errors introduced in this document, and
claims no rights to any of the mathematical content.}

\tableofcontents

\providecommand{\scr}[1]{{\mathscr{#1}}}
\renewcommand{\cal}[1]{{\mathcal{#1}}}
\renewcommand{\frak}[1]{{\mathfrak{#1}}}
\renewcommand{\geq}{\geqslant}
\renewcommand{\leq}{\leqslant}

\providecommand{\sh}[1]{\scr{#1}}

\providecommand{\OO}{\sh{O}}
\providecommand{\HH}{\mathrm{H}}
\providecommand{\cc}{\mathrm{c}}
\providecommand{\coh}{\mathrm{coh}}
\providecommand{\CC}{\mathrm{C}}
\providecommand{\RR}{\mathrm{R}}
\providecommand{\LL}{\mathrm{L}}
\providecommand{\TT}{\mathbf{T}}
\providecommand{\KK}{\mathbf{K}}
\providecommand{\DD}{\mathrm{D}}

\providecommand{\Hom}{\operatorname{Hom}}
\providecommand{\Ext}{\operatorname{Ext}}
\providecommand{\shExt}{\mathscr{E}\mathcal{xt}}
\providecommand{\Supp}{\operatorname{Supp}}

\section*{Introduction}\label{introduction}
\addcontentsline{toc}{section}{Introduction}

\oldpage{261}

Given a \(\sigma\)-compact analytic space \(X\), of bounded dimension,
we know how to construct \autocite{10} a bounded complex
\(\mathbf{K}_X^\bullet\) of \({\mathscr{O}}_X\)-modules with coherent
cohomology, and a (continuous linear) ``trace''
\(\mathbf{T}_X\colon\mathrm{H}_\mathrm{c}^0(X;\mathbf{K}_X^\bullet)\to\mathbb{C}\)
that can happily provide a duality between the \textbf{TODO:
separations?} of the spaces \(\mathrm{H}_\mathrm{c}^p(X;{\mathscr{F}})\)
and \(\operatorname{Ext}^{-p}(X;{\mathscr{F}},\mathbf{K}_X^\bullet)\),
for every integer \(p\) and every coherent \({\mathscr{O}}_X\)-module
\({\mathscr{F}}\), as long as we endow these space with suitable
topologies; similarly, the trace induces a perfect pairing between
\(\mathrm{H}^p(X;{\mathscr{F}})\) and
\(\operatorname{Ext}_\mathrm{c}^{-p}(X;{\mathscr{F}},\mathbf{K}_X^\bullet)\),
suitably topologised. The purpose of this article is the study of the
compatibility of this duality with direct images: if \(f\colon X\to Y\)
is a morphism of analytic spaces, then we have to construct a ``relative
trace''
\(\mathbf{T}_{X/Y}\in\operatorname{Ext}^0(Y;\mathrm{R}f_!\mathbf{K}_X^\bullet,\mathbf{K}_Y^\bullet)\)
(where \(f_!\) denotes the direct image with proper support). We then
have, for every complex \({\mathscr{F}}^\bullet\) of
\({\mathscr{O}}_X\)-modules with coherent cohomology
(\({\mathscr{F}}^\bullet\in\mathrm{D}_\mathrm{coh}(X)\)) and for every
complex \({\mathscr{G}}^\bullet\) with bounded-below coherent cohomology
(\({\mathscr{G}}^\bullet\in\mathrm{D}_\mathrm{coh}^+(Y)\)), a canonical
morphism \[
  \begin{aligned}
    \mathrm{R}f_!{\mathscr{F}}\operatorname{Hom}\big(
      X ; {\mathscr{F}}^\bullet,
      \mathrm{R}\operatorname{Hom}(
        \mathrm{L}f^*&\mathrm{R}\operatorname{Hom}(
          Y;{\mathscr{G}}^\bullet,
          \mathbf{K}_Y^\bullet
        ),
        \mathbf{K}_X^\bullet)
    \big)
  \\\to\, &\mathrm{R}\operatorname{Hom}(Y;\mathrm{R}f_*{\mathscr{F}}^\bullet,{\mathscr{G}}^\bullet)
  \end{aligned}
\tag{1}
\] that is, in certain cases, an isomorphism (we obtain another formula
by swapping \(\mathrm{R}f_!\) and \(\mathrm{R}f_*\), and more striking
formulas by setting \({\mathscr{G}}^\bullet=\mathbf{K}_Y^\bullet\)).

In general, direct images are not coherent, and, just like the absolute
case, we can only hope to address these questions by defining a
topological dualisation
\(\mathrm{R}\operatorname{Hom}(Y;-,\mathbf{K}_Y^\bullet)\) on \(Y\) (the
relative trace should live in a topological \(\operatorname{Ext}^0\),
and (1) only gives isomorphisms of suitable operations up to
``separation''). We restrict ourselves here to the case where \(f\) is
proper, and we use the Grauert direct image theorem.\footnote{In another
  article, J.P.~Ramis and G.~Ruget hope to treat the general case, once
  there is a generalisation to the relative case of the theorems of
  Andreotti--Grauret on \(p\)-convex spaces. The above theorem would
  then follow from the conjunction of the general duality theorem and
  the Grauert direct image theorem.} In this case, (1) is an isomorphism
(the \(\operatorname{Hom}\) being the usual functors), and we have an
``algebraic'' statement. \oldpage{262} The proof, however, uses
topology: to establish the existence of the relative trace, and then to
prove that (1) is an isomorphism; we invoke, in particular, the
separation of certain cohomology spaces.

\hyperref[section-4]{Section 4}, which deals with questions of
separation, is independent of the previous sections. There we show the
two following theorems, conjectured by Malgrange, and which form an
indispensable complement to the absolute duality theorems:

\begin{itenv}{Theorem 4.1}
Given an analytic space \(X\), a coherent sheaf \({\mathscr{F}}\), and
an integer \(p\), suppose that every compact \(K\subset X\) has a
compact ``envelope'' in \(K'\) such that every element of
\(\mathrm{H}_K^p(X;{\mathscr{F}})\) whose image in
\(\mathrm{H}_\mathrm{c}^p(X;{\mathscr{F}})\) is zero already has zero
image in \(\mathrm{H}_{K'}^p(X;{\mathscr{F}})\). Then the space
\(\mathrm{H}_\mathrm{c}^p(X;{\mathscr{F}})\), endowed with the topology
constructed in \autocite{10}~(the only quotient topology of the
Fréchet--Schwartz dual compatible with the duality theorem), is
separated.

\end{itenv}

As for the statement of Theorem~4.9, we obtain this by replacing, in the
lines above, \(\mathrm{H}_\bullet^p(X;{\mathscr{F}})\) with
\(\operatorname{Ext}_\bullet^p(X;{\mathscr{F}},\mathbf{K}_X^\bullet)\).

In an appendix, we will give another proof, using Theorem~4.9, of a
theorem of Siu on the vanishing of certain cohomology groups.

\section{Sketch of a proof}\label{section-1}

We denote by \(\iota\) the canonical isomorphism from
\(\mathrm{H}_\mathrm{c}^0(Y;\mathrm{R}f_*\mathbf{K}_X^\bullet)\) to
\(\mathrm{H}_c^0(X;\mathbf{K}_X^\bullet)\). The trace \(\mathbf{T}_Y\)
induces a pairing between
\(\mathrm{H}_\mathrm{c}^0(Y;\mathrm{R}f_*\mathbf{K}_X^\bullet)\) and
\(\operatorname{Ext}^0(Y;\mathrm{R}f_*\mathbf{K}_X^\bullet,\mathbf{K}_Y^\bullet)\),
and we are looking for an element \(\mathbf{T}_{X/Y}\) in this latter
space such that we have, for all
\(a\in\mathrm{H}_\mathrm{c}^0(Y;\mathrm{R}f_*\mathbf{K}_X^\bullet)\), \[
  \mathbf{T}_Y(a\cdot\mathbf{T}_{X/Y})
  = \mathbf{T}_x\iota a.
\] In other words, we are looking for a transpose \({}^t\iota\) of
\(\iota\) with respect to the pairings \(\mathbf{T}_Y\) and
\(\mathbf{T}_X\) (with \(\mathbf{T}_{X/Y}\) as the image of
\({}^t\iota\) under the function
\(1\in\mathrm{H}^0(X;{\mathscr{O}}_X)\)). Since this problem is local on
\(Y\) (provided that we are guaranteed the local vanishing of
\(\operatorname{Ext}^i(Y;\mathrm{R}f_*\mathbf{K}_X^\bullet,\mathbf{K}_Y^\bullet)\)
for \(i<0\)), we can suppose that \(Y\) is realisable as an analytic
subset of an open polydisc \(V\) of dimension~\(p\). If we still denote
by \(f\) the composite of \(f\) and an embedding of \(Y\) into \(V\), it
follows from the construction of \(\mathbf{K}_Y^\bullet\) that what we
are looking for is an element of
\(\operatorname{Ext}^p(V;\mathrm{R}f_*,\mathbf{K}_X^\bullet,\Omega_V)\).
We now want to proceed locally on \(X\), supposing (being content to
shrink \(X\)) that \(X\) is the union of a finite number of opens
\(X_i\), each of which is embeddable into an open \(U_i\) of
\(\mathbb{C}^n\): we can incorporate each \(X_i\) into a commutative
diagram

\pandocbounded{\includegraphics[keepaspectratio]{IM-13-1971-261_files/mediabag/_svgs/1a05921ba0803713f5ce69d9b39322e4afcab572.pdf}}

\oldpage{263}

where \(I\) is an embedding, and where \(F\) denotes the second
projection of \(U_i\times V\). We denote by \(f_{i_0i_1\ldots i_r}\) the
direct image with proper support in
\(X_{i_0}\cap X_{i_1}\cap\ldots X_{i_r}\) and we first search for a
``trace'' \(\mathbf{T}_i\) in each
\(\operatorname{Ext}^p(V;\mathrm{R}f_i\mathbf{K}_X^\bullet,\Omega_V)\).


% Bibliography
\nocite{*}
\printbibliography[heading=bibintoc,title=Bibliography]

\end{document}
